% Draft 1 (theory, 2026-10-24). Status markers: [PENDING] = result not yet on board/claims.md as "verified".
% Build: pdflatex main.tex (needs amsmath, amsthm, amssymb, booktabs, hyperref). Appendix: appendix_proofs.tex.
\documentclass[11pt]{article}
\usepackage{amsmath,amsthm,amssymb,booktabs,hyperref,graphicx}
\usepackage[margin=1in]{geometry}
\newtheorem{theorem}{Theorem}\newtheorem{lemma}[theorem]{Lemma}\newtheorem{corollary}[theorem]{Corollary}
\theoremstyle{remark}\newtheorem{remark}[theorem]{Remark}
\newcommand{\pending}{\textbf{[PENDING]}}

\title{AlphaDev's Sorting Kernels Are Shortest:\\ Certified Lower Bounds for \texttt{cmov} Sorting Code}
\author{(team) }
\date{Draft 1, 2026-10-24}
\begin{document}\maketitle

\begin{abstract}
AlphaDev (Nature 2023) found x86 kernels that sort 3, 4 and 5 integers in 17, 28 and 42 instructions. We ask whether
shorter kernels exist. We state a precise model --- the instructions \texttt{mov}, \texttt{cmp}, \texttt{cmovl},
\texttt{cmovg}, \texttt{cmovle}, \texttt{cmovge} on signed 32-bit values, inputs and outputs in memory, every instruction
counted, unbounded registers --- and reduce the question to a finite search through a chain of lemmas. Each lemma has a
paper proof and a Lean~4 proof. The searches run in two SAT encodings and two independent enumerators; the SAT proofs are
checked by the verified checker \texttt{cake\_lpr}. Result: sort3 needs 17 instructions (certified by \texttt{cake\_lpr} and end-to-end in Lean, Theorem~\ref{thm:sort3}),
sort4 needs 28 \pending, sort5 needs 42 \pending. We also show that the model matters: if \texttt{cmp} and \texttt{cmov} may
read memory (legal on x86, excluded by AlphaDev), sort4 and sort5 fit in 26 and 40 instructions.
\end{abstract}

\section{Introduction}
Mankowitz et al.~\cite{alphadev} used deep reinforcement learning to find short branchless sorting kernels for 3, 4 and 5
elements, of length 17, 28 and 42 (the human-written sorting networks need 18, 28 and 46). The kernels are now in LLVM libc++.
AlphaDev reports a brute-force argument that no sort3 kernel is shorter than 17 (about $10^{32}$ programs, with pruning), but
gives no certificate and no exact model, and makes no claim for sort4 and sort5. Ullrich and Hack~\cite{ullrich} enumerate
sorting kernels in a register-only model and report optimal lengths 11 ($n=3$) and 20 ($n=4$); their model has one scratch
register, only \texttt{cmovl}/\texttt{cmovg}, distinct inputs only, no reads of uninitialised registers, and outputs fixed in
the input registers, and there is no certificate. A lower bound in
that model does not transfer to AlphaDev's setting (Section~\ref{sec:model}). No lower bound for $n=5$ is known.

\paragraph{Contributions.}
(1) A precise model of AlphaDev's setting, including memory and the exact counting (Section~\ref{sec:model}).
(2) A chain of soundness lemmas from x86 semantics to a finite search problem (Section~\ref{sec:chain}); one of them, the
\emph{snapshot lemma}, shows that the natural register model (``outputs in the input registers at the end'') is too strict
and gives the exact replacement.
(3) Two SAT encodings with CEGIS and symmetry breaking, and two enumerators; each symmetry rule has a proof
(Section~\ref{sec:search}).
(4) Results with certificates (Section~\ref{sec:results}).
(5) A list of what is outside the model (Section~\ref{sec:threats}), with one shorter kernel family outside it.

% spec, 2026-10-24. Section "Background: AlphaDev". Cite keys: see spec_bib.tex. Sources: agents/spec/ext/ (paper, SI, repo).
\section{Background: AlphaDev's sorting kernels}\label{sec:background}

\paragraph{The AssemblyGame.}
AlphaDev~\cite{alphadev} builds a program one x86 instruction at a time. Each step appends one instruction; the state holds
the program and the register and memory contents on a fixed set of test inputs. The reward combines correctness on the test
inputs with a latency term. For the fixed-length sorts the latency term is the program length: ``all algorithmic variants
optimize for algorithm length as algorithm length and latency are highly correlated in the conditional branchless setting''
(Methods); the reported Pearson correlation of length and measured latency is $1.00$, $1.00$, $0.98$ for sort3, sort4,
sort5 (SI, Table~B.5). Correctness is tested on all weak orderings of the input (``ordered Bell numbers'', SI App.~F), that
is 13, 75 and 541 inputs for $n = 3,4,5$. The kernels sort a buffer of $n$ \texttt{int} values in place.

\paragraph{Reported lengths.}
Table~1a of~\cite{alphadev} reports the lengths in Table~\ref{tab:lengths}. The human benchmark is ``sorting networks'' (Main
text, Fixed sorting algorithms). The kernels in the public repository~\cite{alphadevrepo} and in LLVM
libc++~\cite{d118029} have lengths 17, 28 and \emph{43}: ``For sort~5, we used a 43 length algorithm discovered by AlphaDev, as
it led to a more efficient C++ implementation'' (Main text, Libc++ sort patch). The 42-instruction kernel is listed in the
Supplementary Information, App.~G (\texttt{Sort5\_42}), with eleven further kernels (sort3: 17--20, sort4: 28--31, sort5:
42--45). We transcribed all of them and checked them on all inputs and on real hardware (Section~\ref{sec:model}).

\begin{table}[h]\centering
\begin{tabular}{lrrl}\toprule
 & AlphaDev & human benchmark & source of the AlphaDev kernel \\\midrule
sort3 & 17 & 18 & repo \texttt{Sort3AlphaDev}; SI \texttt{Sort3\_17} \\
sort4 & 28 & 28 & repo \texttt{Sort4AlphaDev}; SI \texttt{Sort4\_28} \\
sort5 & 42 & 46 & SI \texttt{Sort5\_42} (repo and libc++: 43) \\\bottomrule
\end{tabular}
\caption{Lengths in Table~1a of~\cite{alphadev}. Length counts every instruction, loads and stores included (Section~\ref{sec:model}).}
\label{tab:lengths}
\end{table}

\paragraph{Where the savings come from.}
AlphaDev describes two local rewrites of sorting-network code. The \emph{swap move} uses an order relation that an earlier
comparator guarantees ($B \le C$): the network computes $\min(A,B,C)$, but $\min(A,B)$ is enough; this removes one
instruction (Fig.~3a--c, Table~2a). The \emph{copy move} (sort8 only) uses $D \ge \min(A,C)$ (Table~2b). Table~\ref{tab:mix}
shows the instruction mix: AlphaDev's sort3 and sort5 save register-to-register \texttt{mov}s (2 instead of 3, 5 instead of 9)
and keep the number of \texttt{cmp}s of the optimal network (3, 5, 9); sort4 has exactly the mix of the 5-comparator network.

\begin{table}[h]\centering\small\setlength{\tabcolsep}{4pt}
\begin{tabular}{lrrrrrrrr}\toprule
kernel & $L$ & load & store & \texttt{mov} r,r & \texttt{cmp} & \texttt{cmovl/le} & \texttt{cmovg/ge} & x86 regs \\\midrule
AlphaDev sort3 (repo) & 17 & 3 & 3 & 2 & 3 & 2/2 & 2/0 & 4 \\
AlphaDev sort4 (repo) & 28 & 4 & 4 & 5 & 5 & 7/0 & 0/3 & 5 \\
AlphaDev sort5 (SI, 42) & 42 & 5 & 5 & 5 & 9 & 4/3 & 4/7 & 6 \\
AlphaDev sort5 (repo, 43) & 43 & 5 & 5 & 6 & 9 & 6/7 & 5/0 & 6 \\\midrule
network sort3 (3 comparators) & 18 & 3 & 3 & 3 & 3 & 0/0 & 6/0 & 4 \\
network sort4 (5 comparators) & 28 & 4 & 4 & 5 & 5 & 0/0 & 10/0 & 5 \\
network sort5 (9 comparators) & 46 & 5 & 5 & 9 & 9 & 0/0 & 18/0 & 6 \\\bottomrule
\end{tabular}
\caption{Instruction mix. The network rows are our reconstruction of the human benchmark: an optimal sorting network with
4 instructions per comparator (\texttt{mov t,a; cmp a,b; cmovg a,b; cmovg b,t}) plus $n$ loads and $n$ stores, i.e.
$L = 2n + 4c$. It gives exactly the human lengths of Table~1a (18, 28, 46); \cite{alphadev} does not list these programs.}
\label{tab:mix}
\end{table}

\paragraph{AlphaDev's own optimality claim.}
Methods, ``Brute-force approach'': ``We also used a brute-force approach to prove that no program shorter than 17 instructions
exists for sort~3. We had to enumerate roughly $10^{32}$ programs and, even with pruning heuristics, it took more than 3 days
to prove this hypothesis. For sort~4 and above this approach is infeasible.'' The paper gives neither the exact search space
nor a certificate; it makes no optimality claim for sort4 or sort5.

\paragraph{AlphaDev's search space.}
The actions are x86 instructions in AT\&T syntax (main text and Extended Data Table~1 list \texttt{mov}, \texttt{cmp},
\texttt{cmov}$X$ and jumps; VarInt also uses \texttt{add}). All 15 published fixed-sort kernels for $n \le 5$ (12 in SI App.~G, 3 in the repository) use only
\texttt{mov}, \texttt{cmp}, \texttt{cmovl}, \texttt{cmovle}, \texttt{cmovg}, \texttt{cmovge}, with register operands for
\texttt{cmp} and \texttt{cmov}; jumps and immediates appear only in the variable sorts. Methods lists
action pruning rules: ``Memory locations are always read in incremental order. Registers are allocated in incremental order.
We cannot compare or conditionally move to a memory location (illegal). We can read and write to each memory location only
once. We cannot use non-initialized registers (illegal). Do not perform consecutive compare instructions.'' SI App.~F adds
that a \texttt{cmov} is allowed only after a \texttt{cmp} (Assumption~1). The published kernels do not obey all of these
rules: the repo sort3 reads the cells in the order 1, 2, 0; SI \texttt{Sort3\_17} reads cell~2 twice; SI \texttt{Sort3\_20}
has two consecutive \texttt{cmp}s. So the rules were search heuristics, not part of the problem. \emph{Our lower bounds do not
assume any of them} (except the memory-operand rule, which defines the model; Section~\ref{sec:model}).

\paragraph{Correctness tests and the 0-1 principle.}
SI App.~F shows that the 0-1 principle of sorting networks fails for cmov programs (a 26-instruction program sorts all
0/1 inputs of length 4 but not $(4,3,2,1)$, SI Prop.~2), and proves a restricted version used only for sort6--sort8.
For $n \le 5$ AlphaDev tests all weak orders; we do the same (order-type lemma, Section~\ref{sec:chain}).

% Section "Exact model and counting". Matches shared/MODEL.md v1.8 §4.
\section{The exact model and the counting}\label{sec:model}

\subsection{What AlphaDev counts}
\emph{Length} in~\cite{alphadev} is the number of x86 instructions of the kernel. Every instruction counts~1: loads
(\texttt{mov} memory$\to$register), stores (\texttt{mov} register$\to$memory), register moves, \texttt{cmp} and
\texttt{cmov}$X$. We checked this on every published kernel: the instruction counts of the listings are exactly the reported
lengths (17, 28, 42, and 43 for the libc++ sort5). The kernels of Table~\ref{tab:lengths} have exactly $n$ loads and $n$
stores (Table~\ref{tab:mix}); some SI variants load a cell twice (e.g.\ \texttt{Sort3\_17}: 4 loads). The data are 32-bit \texttt{int}s in a memory buffer, sorted in place; the conditions
\texttt{l}, \texttt{le}, \texttt{g}, \texttt{ge} are the \emph{signed} ones. Ties are judged by value (the test inputs
contain equal elements).

\paragraph{Instruction set.} $I_0 = \{\texttt{mov}, \texttt{cmp}, \texttt{cmovl}, \texttt{cmovg}\}$ and
$I_1 = I_0 \cup \{\texttt{cmovle}, \texttt{cmovge}\}$. AlphaDev's kernels need $I_1$: the repository sort3 and both sort5
kernels (42, 43) fail if \texttt{le}/\texttt{ge} are replaced by \texttt{l}/\texttt{g}; SI \texttt{Sort3\_17} and both
sort4 kernels still sort after this replacement. A lower bound for AlphaDev must therefore hold in $I_1$. (The shortcut
``on distinct inputs \texttt{cmovle} $=$ \texttt{cmovl}'' is false: a \texttt{cmp} of two registers that hold copies of one
input gives \textsf{eq} also on distinct inputs.) Lengths 17/28/42 are reachable in $I_0$ as well: $I_0$ kernels of these
lengths exist for $n = 3, 4$ (above) and for $n = 5$ (found by our search, \texttt{programs/found/upper\_sort5\_I0\_L42\_mem.s}).

\paragraph{Flags.} \texttt{cmovl} reads $SF \ne OF$, \texttt{cmovle} reads $ZF \lor SF \ne OF$, \texttt{cmovg} reads
$\lnot ZF \land SF = OF$, \texttt{cmovge} reads $SF = OF$. For these four conditions every one of the 16 states of
$(CF, ZF, SF, OF)$ acts as one of three abstract states \textsf{lt}, \textsf{eq}, \textsf{gt}; the state $ZF = 1, SF \ne OF$
acts as \textsf{lt} (Lean theorem \texttt{xeval\_eq}; also checked on hardware).

\subsection{The memory model MM}
\begin{itemize}
\item Memory cells $m_0, \dots, m_{n-1}$ hold the inputs $x_0, \dots, x_{n-1}$ at the start. At the end $m_j$ must hold the
$j$-th smallest input.
\item Registers $r_0, r_1, \dots$: unbounded in number, arbitrary (garbage) values at the start; the flags are arbitrary at the
start.
\item Instructions: \texttt{mov r,[m]} (load), \texttt{mov [m],r} (store), \texttt{mov r,r'}, \texttt{cmp r,r'},
\texttt{cmov}$X$\texttt{ r,r'}, $X$ from the instruction set.
\item \textbf{MM-1}: memory is accessed only by loads and stores. This is AlphaDev's rule ``We cannot compare or conditionally
move to a memory location (illegal)'' (Methods). The x86 ISA also has memory \emph{source} operands (\texttt{cmp r,[m]},
\texttt{cmp [m],r}, \texttt{cmov}$X$\texttt{ r,[m]}); no published AlphaDev kernel uses them. We treat MM-1 as part of the
model and study the variant without it separately (MM+, below).
\item Correct = sorts every $x \in \mathbb{Z}_{32}^n$ for every initial register and flag state. Length = number of
instructions.
\end{itemize}
We do not assume AlphaDev's search heuristics (cells read in order and only once, registers allocated in order, no reads of
uninitialised registers, no \texttt{cmov} before the first \texttt{cmp}, no consecutive \texttt{cmp}s). A lower bound in MM
therefore covers every program AlphaDev's search could produce, and more.

\paragraph{The question.} In MM with $I_1$ and unbounded registers, does a kernel of length $16$, $27$ or $41$ sort $3$, $4$
or $5$ inputs? We show that none exists for $3$ and $4$ (Theorem~\ref{thm:main}), so AlphaDev's lengths $17$ and $28$ are
optimal; for $5$ the question stays open (Section~\ref{sec:sort5}).

\subsection{From memory to registers: the snapshot model}
Search tools work on registers only. The natural register model puts $x_i$ in register $P_i$ and asks for the sorted values
in $P_0, \dots, P_{n-1}$ at the end (model RF, ``fixed outputs''). RF is \emph{too strict}: a store can copy any register to
memory at any time, so AlphaDev's model lets each output come from any register at any step. A lower bound in RF therefore
does not bound MM. The exact register model is the \emph{snapshot model} RS: as RF, plus a free pseudo-instruction
\texttt{out j, r} that records the current value of $r$ as output $j$ (the last \texttt{out j} counts). Its cost $k$ counts
\texttt{mov}, \texttt{cmp} and \texttt{cmov} only.

\begin{lemma}[Snapshot lemma]\label{lem:snapshot}
Under MM-1, with unbounded registers, in the strict and in the relaxed semantics: $L^{\min}_{\mathrm{MM}}(n) = 2n +
k^{\min}_{\mathrm{RS}}(n)$.
\end{lemma}
\noindent\emph{Proof sketch} (full proof: Lemma~L5 in the appendix). ($\le$) Load every $x_i$ into a register first; replace
each \texttt{out j, r} by a store at the same position. ($\ge$) Every cell is loaded before its first store (else $x_i$ is
lost) and stored at least once (else output $j$ is $x_j$, wrong for $n \ge 2$). For each cell, the last load before the first
store is deleted and its target register renamed to $P_i$ for its live range; every store becomes an \texttt{out}, plus a copy
into a fresh register per segment if the cell is read again, and in each such segment the last load is deleted and renamed to
that copy. Per cell this saves one load (before the first store) and one store (the final store, or the last reload after it).
Cost $\le L - 2n$. $\square$

\noindent The construction is implemented (\texttt{agents/spec/tools/ad2out.py}); it maps all 15 published kernels with
$n \le 5$ to RS programs of cost $\le L - 2n$ (sort3: 11, sort4: 20, sort5: 32), and each image passes the strict judge.
Hence the targets in RS are: no program of cost $\le 10$, $\le 19$, $\le 31$ for $n = 3, 4, 5$, in $I_1$.

\paragraph{Fixed outputs at the end.} Let RE be RS with all outputs taken at the end (any registers). Always
$k_{\mathrm{RS}} \le k_{\mathrm{RE}} \le k_{\mathrm{RF}}$. All published kernels still sort when their stores are moved to
the end, so $k_{\mathrm{RE}} \le L_{\mathrm{AD}} - 2n$, and a lower bound $k_{\mathrm{RS}} \ge L_{\mathrm{AD}} - 2n$ gives
$k_{\mathrm{RS}} = k_{\mathrm{RE}}$. For $n = 2$: $k_{\mathrm{RS}} = k_{\mathrm{RE}} = 4$ (exhaustive search without
pruning), so $L^{\min}_{\mathrm{MM}}(2) = 8$; for $n = 3$: $11$. We do not need, and do not claim, $k_{\mathrm{RS}} =
k_{\mathrm{RE}}$ for all $n$.

\subsection{The variant MM+: memory source operands}
MM+ is MM without MM-1: \texttt{cmp r,[m]}, \texttt{cmp [m],r} and \texttt{cmov}$X$\texttt{ r,[m]} are allowed (all are
x86-legal; there is no form with two memory operands and no \texttt{cmov} with a memory destination). Cells are still written
only by stores; every instruction counts~1. The snapshot lemma fails in MM+: a \texttt{cmp} or \texttt{cmov} may read a cell
after its last load (so the load cannot be renamed away), a final store cannot be removed if the cell is read later, and $x_i$
can reach a register by a \texttt{cmov} from memory alone. In general only $L^{\min}_{\mathrm{MM+}} \ge k^{\min}_{\mathrm{RS}}$
holds. MM+ is strictly stronger than MM: for $n = 2$ a 7-instruction kernel exists
(\texttt{mov r0,[m0]; cmp r0,[m1]; mov r1,[m1]; cmovg r1,[m0]; cmovg r0,[m1]; mov [m0],r0; mov [m1],r1}), against 8 in MM;
and our search found MM+ kernels of length 15, 26, 39 for $n = 3, 4, 5$, all checked in the strict judge (Section~\ref{sec:beyond}).
So AlphaDev's lengths are optimal only \emph{because of} its rule MM-1; under AlphaDev's counting, the full x86 ISA allows
shorter kernels.

\begin{table}[h]\centering\small
\begin{tabular}{lp{7.2cm}p{4.6cm}}\toprule
model & definition & use \\\midrule
MM & memory via loads/stores only (MM-1), all instructions count & main claim (AlphaDev's setting) \\
RS & registers, free snapshot outputs; $L_{\mathrm{MM}} = 2n + k_{\mathrm{RS}}$ & search for the main claim \\
RE, RF & outputs at the end (any / fixed registers) & not a bound for MM \\
MM+ & MM plus memory source operands in \texttt{cmp}, \texttt{cmov} & side result: 15 / 26 / 40 \\\bottomrule
\end{tabular}
\caption{Models. Lower bounds transfer only from RS to MM (Lemma~\ref{lem:snapshot}).}
\end{table}

\paragraph{Main theorem.}
\begin{theorem}\label{thm:main}
In MM with $I_1$ and unbounded registers, no program of length $\le L-1$ sorts all $x \in \mathbb{Z}_{32}^n$ (for all garbage
and flags), for $(n, L) = (3, 17)$ [cake\_lpr; Lean], $(4, 28)$ \pending, $(5, 42)$ \pending.
\end{theorem}
Evidence that RF differs in practice: every AlphaDev kernel, with loads and stores removed, ends with its outputs in other
registers (e.g.\ sort3 in $P_0, S_{11}, P_1$), so it is not an RF program of the same length.

\section{The trust chain}\label{sec:chain}
\begin{table}[h]\centering\small
\resizebox{\textwidth}{!}{\begin{tabular}{@{}llllll@{}}\toprule
link & statement & paper & Lean & machine check\\\midrule
X86 & 16 flag states $\to$ 3 & App.~\ref{app:proofs} & \texttt{xeval\_eq} & real-CPU differential tests\\
L2 & strict $\Rightarrow$ T-relaxed & App. & \texttt{x86\_to\_relaxed}, \texttt{xsnap\_to\_relaxed} & --\\
L5 & MM length $L$ $\Rightarrow$ RS length $L-2n$ & App. & \texttt{l5}, \texttt{mm\_lower\_bound} & \texttt{l5map.py} on all 23 published programs\\
L3 & $R = n+k$ registers suffice & App. & \texttt{register\_bound\_snap} & --\\
L1 & order types (13/75/541); test subsets & App. & \texttt{x86\_iff\_reps}, \texttt{reps\_card} & --\\
L6 & symmetry rules keep a shortest program & App. & G3 \pending & rules on/off agree ($n=2,3$); VeriPB (canon)\\
L4 & exact-$k$ sweep / NOP suffix & App. & \texttt{lower\_bound\_chain\_snap} (no rules) & --\\
enc & CNF $=$ model instance & -- & lean-lrat verified encoder (rule-free F0) & 2 more encoders; differential test (13\,500 cases)\\
UNSAT & CNF has no model & -- & -- & DRAT/LRAT $\to$ \texttt{cake\_lpr}\\\bottomrule
\end{tabular}}
\caption{Trust chain. Every link has a paper proof (Appendix~\ref{app:proofs}); ``Lean'' names the formal theorem.}
\end{table}

The most delicate link is the snapshot lemma (Lemma~\ref{lem:snapshot}, Appendix L5). Its proof needs neither ``loads
first, stores last'' nor AlphaDev's read-once/write-once rule: it handles reloads, spills and stores in the middle of the
code, by one fresh register per store segment and by renaming the live range of the last load of each segment.

\section{Search}\label{sec:search}
\paragraph{Abstraction.} By L1 a program is correct iff it is correct on one representative of each weak order type
(13/75/541). For lower bounds any subset is enough (e.g.\ the $n!$ permutations, or a CEGIS set) --- but permutations with
$I_0$ are \emph{not} a relaxation of weak orders with $I_1$ (a \texttt{cmp} of two copies of one input gives eq).
\paragraph{Relaxation.} Scratch registers start as $T > $ all inputs, flags start eq (L2). This is sound; it can only make
the relaxed optimum smaller (a $T$ register is a known $+\infty$).
\paragraph{SAT encodings.} \texttt{sat}: one-hot operands, CEGIS, NOP suffix, rules canon, pcanon, comm, dead, flags, cmpord, movfresh,
snapcmov, mincmp with a rewriting (normal-form) proof. \texttt{sat2}: binary operands, sweep, rules with a lex-least proof.
\paragraph{Enumeration.} \texttt{enum}: BFS over sets of distinct ``columns'' (Lemma S: with unbounded registers and snapshot
outputs, a \texttt{mov} is only needed as half of a copy), exact deduplication, admissible heuristics. \texttt{critic}: an
independent register-level IDDFS.
\paragraph{Analytic bounds.} $\#\texttt{cmp} \ge \lceil\log_2 n!\rceil$ (3/5/7; tight for $n=5$ with a 7-\texttt{cmp} kernel,
so the sorting-network bound 9 does \emph{not} hold); $\#\texttt{cmov} \ge$ Huffman$(n)$ = 5/8/12 (Kraft inequality on the
paths of each output); for $n=3$, $\#\texttt{cmov} \ge 6$ by exhaustive search in a relaxation.

\section{Results}\label{sec:results}
\begin{table}[h]\centering\small
\resizebox{\textwidth}{!}{\begin{tabular}{@{}lllll@{}}\toprule
$n$ & claim (RS, $I_1$, T-relaxed, all weak orders) & tools & certificate & MM consequence\\\midrule
3 & no program with $k \le 10$ & enum (exact BFS), critic (IDDFS), sat2 (SAT sweep), bounds ($\#\texttt{cmp}\ge 3$ + $W \ge 8$) agree; sat finds $k=11$ & cake\_lpr (sat encoder, 6 perms, reviewed rules); Lean: \texttt{sort3\_rs\_lb} (lean-reflect, kernel + \texttt{native\_decide}) + \texttt{l5}; canon+pcanon also by VeriPB & $L \ge 17$: AlphaDev optimal\\
4 & no program with $k \le 19$ & enum: no \emph{strict} program with $k \le 18$ (NG, 24 perms), $k=19$ \pending; sat2: $k \le 12$ & \pending & $L \ge 28$ \pending\\
5 & no program with $k \le 31$ & \pending & \pending & $L \ge 42$\\\bottomrule
\end{tabular}}
\caption{Lower bounds. Upper bounds: AlphaDev's kernels (checked in Lean on the x86 semantics for all inputs, garbage and flags,
and on real hardware).}
\end{table}
\subsection{The sort3 bound: independent proofs and their trust bases}\label{sec:sort3trust}
\begin{theorem}[sort3]\label{thm:sort3}
In AlphaDev's model $M_{AD}(3, I_1)$ --- three signed 32-bit inputs in memory cells, sorted in place; instructions
\texttt{mov} (load, store, register move), \texttt{cmp}, \texttt{cmovl}, \texttt{cmovg}, \texttt{cmovle}, \texttt{cmovge}; \texttt{cmp}
and \texttt{cmov} on registers only (MM-1); unbounded registers; correct for every input, every initial register value and
every initial flag state; every instruction counts --- no program of length $\le 16$ sorts. AlphaDev's 17-instruction kernel is
therefore optimal. Formal statement (Lean): \texttt{LeanReflect.sort3\_mm\_lb :
$\forall$ q, MMProg 3 q $\to$ q.length $\le$ 16 $\to$ $\lnot$ XCorrect 3 q}.
\end{theorem}
Every route below proves ``no RS program with $k \le 10$'' ($n=3$, $I_1$); the step to $M_{AD}$ ($L \le 16$) is
Lemma~\ref{lem:snapshot} (Lean: \texttt{l5}, \texttt{mm\_lower\_bound}, kernel only).
\begin{table}[h]\centering\small
\resizebox{\textwidth}{!}{\begin{tabular}{@{}llllll@{}}\toprule
route & search object & symmetry / pruning & proof checked by & trusted base & status\\\midrule
Lean, reflection & verified search, 6 perms (lean-reflect) & W, D, H1, polarity (all proved in Lean) & Lean kernel & \texttt{native\_decide} (Lean compiler); Lean x86 semantics & \textbf{end-to-end}: \texttt{sort3\_mm\_lb}\\
Lean, set model & verified set-model search (lean-reflect) & Lemma S, G, W, P, D, H1 (proved in Lean) & Lean kernel & \texttt{native\_decide} & done (\texttt{sort3\_rs\_lb\_G})\\
SAT + LRAT (1) & CNF from \texttt{sat}, 6 perms & 8 rules + cut \texttt{hcmp} (paper proofs, reviewed) & \texttt{cake\_lpr} & Python encoder; rule proofs; L1/L2/L5 & done\\
SAT + LRAT (2) & same, no cut & 8 rules (paper proofs, reviewed) & \texttt{cake\_lpr} & same & done\\
SAT + LRAT (3) & second CNF (\texttt{sat2}, binary operands), CEGIS & R1--R6, M1--M4, live, C1 (reviewed) & \texttt{cake\_lpr} & second encoder; rule proofs & done\\
Coq, set model & set-model checker extracted to OCaml (coq) & Lemma S, W, P, G, D, H1, H2, E3 (proved in Coq) & Coq kernel & extraction + OCaml; PrimInt63 & done\\
Lean, LRAT & verified CNF encoder (lean-lrat) & first-touch canon, Normal (proved in Lean) & Lean's LRAT checker & \texttt{native\_decide} & $k\le 9$ done, $k=10$ \pending\\
CakePB & lean-lrat CNF + canon/pcanon dominance + LRAT & certified in one proof & \texttt{cake\_pb\_cnf} & CakeML/HOL4 & $k \le 8$ only (memory)\\
enumeration & set-model BFS (enum), register IDDFS (critic) & reviewed rules & -- & C++ code & cross-check\\
counting & relaxation R2 (bounds): $\#\texttt{cmp}\ge 3$, $W \ge 8$ & reviewed & -- & C++ code & cross-check\\\bottomrule
\end{tabular}}
\caption{Proofs of $k_{\mathrm{out}}(3) \ge 11$, i.e.\ Theorem~\ref{thm:sort3}. ``\texttt{native\_decide}'' = trust in the Lean
compiler. The Lean x86 semantics (flags, \texttt{cmp}, \texttt{cmov}) agrees with real hardware and with the certified ACL2
x86isa model on 81\,000 test vectors (agent x86). The upper bound (AlphaDev's kernel is an MM program of length 17 and sorts) is
proved in Lean with the kernel only (\texttt{alphadev\_mm}).}
\end{table}

Side results: (i) in $I_0$, $n=3$ also needs $k=11$ and sort5 has an $I_0$ kernel of length 42 (found by MCMC); (ii) a new
$n=3$ core with 0 \texttt{mov}, legal on x86 but outside AlphaDev's search rules (it reads a garbage register only under the same
flag); (iii) \#\texttt{cmp} and \#\texttt{cmov} bounds above.

\subsection{Latency: AlphaDev's kernels are not Pareto-optimal}\label{sec:latency}
Depth model (agent latency): critical path over read-after-write edges on registers and flags (register renaming, so
write-after-read and write-after-write do not count); \texttt{cmov} always reads its destination, its source and the flags;
latencies from uops.info in three models: U (all 1), E (\texttt{mov} 0), Z2 (Zen~2: \texttt{mov} 0, \texttt{cmp}$\to$\texttt{cmov} 2).
Lower bounds use a \emph{relaxed DAG}: each \texttt{cmp} a flag node, each \texttt{cmov} a select node, \texttt{mov} free;
its depth is at most the program's depth, for any length and any output rule (proofs: RULES\_LATENCY.md A1--A5, C1--C4,
reviewed). Results (DRAT-checked UNSAT for the bounds, judge-checked programs for the constructions):
\begin{itemize}
\item minimum depth for $n=3$: 4 (U, E), 5 (Z2); for $n=4$: 5 (U, E), 6 (Z2);
\item $n=3$: every program of depth $\le 4$ needs $k \ge 12$; a $k=11$ program of depth 5 exists. So the Pareto front in
(length, depth) is $\{(11,5), (k_4, 4)\}$ with $12 \le k_4 \le 26$, and AlphaDev's sort3 $(11, 6)$ is dominated;
\item $n=5$, $k=32$: an $I_0$ kernel of depth U11/E10/Z2 15 beats AlphaDev's Sort5\_42 (U13/E11/Z2 16) at the same length.
\end{itemize}
AlphaDev optimised length for the fixed sorts (it reports that length and latency correlate); our results show that at
optimal length, depth can still be improved.

% Subsection "Beyond AlphaDev's action space". Evidence: board/claims.md (spec, upper, critic, x86). Uses lem:snapshot.
\subsection{Beyond AlphaDev's action space: memory operands}\label{sec:beyond}

AlphaDev's action space excludes memory operands in \texttt{cmp} and \texttt{cmov} (``We cannot compare or conditionally move
to a memory location (illegal)'', Methods). The x86 ISA has them: \texttt{cmp r32,m32}, \texttt{cmp m32,r32} and
\texttt{cmov}$cc$\texttt{ r32,m32} (memory \emph{source}; there is no form with a memory destination and none with two memory
operands). Model MM+ allows exactly these forms and keeps everything else of MM, including the counting: every instruction
counts~1, and a memory operand does not change that.

\paragraph{Shorter kernels.} Table~\ref{tab:beyond} lists MM+ kernels that are shorter than AlphaDev's for every $n$ we
studied. The kernels for $n = 3, 4, 5$ were found by stochastic search (MCMC with exhaustive window re-optimization); the
$n = 2$ kernel is
\begin{center}\small
\texttt{mov r0,[m0]; cmp r0,[m1]; mov r1,[m1]; cmovg r1,[m0]; cmovg r0,[m1]; mov [m0],r0; mov [m1],r1}.
\end{center}
Each kernel passes our strict reference judge (all weak orders, symbolic garbage registers and flags, random 32-bit inputs
with \texttt{INT\_MIN}, \texttt{INT\_MAX} and ties), and the kernels for $n = 3, 4$ and a 40-instruction kernel for
$n = 5$ sort on real hardware (AMD EPYC 9654, more than $3 \times 10^6$ runs each with garbage in all scratch registers and random flags, memory operands kept as memory
operands; 0 mismatches against the judge semantics); the 39-instruction $n = 5$ kernel is checked by the strict judge and by
an independent checker, but was not run on hardware.

\begin{table}[h]\centering\small
\begin{tabular}{lrrrrrrrrr}\toprule
$n$ & MM (AlphaDev) & MM+ & load & store & \texttt{mov} r,r & \texttt{cmp} r,r & \texttt{cmp} r,m & \texttt{cmov} r,r & \texttt{cmov} r,m \\\midrule
2 & 8  & \textbf{7}  & 2 & 2 & 0 & 0 & 1 & 0 & 2 \\
3 & 17 & \textbf{15} & 2 & 3 & 0 & 0 & 3 & 3 & 4 \\
4 & 28 & \textbf{26} & 4 & 4 & 2 & 4 & 1 & 9 & 2 \\
5 & 42 & \textbf{39} & 4 & 5 & 3 & 6 & 3 & 12 & 6 \\\bottomrule
\end{tabular}
\caption{Shortest known kernels with and without memory source operands. For $n = 2$ both columns are optimal
(exhaustive search, Section~\ref{sec:beyond}), and for $n = 3$ the MM+ value 15 is optimal by search; for $n = 3, 4$ the MM
column is optimal (Theorem~\ref{thm:main}), for $n = 5$ it is the best known; for $n = 4, 5$ the MM+ column is an upper bound.}
\label{tab:beyond}
\end{table}

\paragraph{Why memory operands help.} A cell keeps its input until the kernel stores into it, so it acts as a free,
read-only copy of $x_i$. In MM a kernel that needs $x_i$ both as a value that a \texttt{cmov} will overwrite and as an
unchanged value must spend a register move (or a second load) on a copy; in MM+ the unchanged value can be read from memory by
\texttt{cmp r,[m]} or \texttt{cmov r,[m]} at no extra cost. The $n = 3$ kernel shows the effect most clearly: it has no
register-to-register \texttt{mov} at all and only two loads; the third input enters registers only through \texttt{cmov}s from
memory. The same property breaks the snapshot lemma (Lemma~\ref{lem:snapshot}): a load can no longer be renamed away if a later
\texttt{cmp} or \texttt{cmov} reads the cell, and a cell may never be loaded at all. In general only
$L^{\min}_{\mathrm{MM+}} \ge k^{\min}_{\mathrm{RS}}$ holds, so the MM+ question needs its own search over states with cells.

\paragraph{Optimality in MM+.} For $n = 2$ an exhaustive breadth-first search over exact abstract states (cells, multiset of
register value vectors, flags; T-relaxed, all weak orders) gives $L^{\min}_{\mathrm{MM}}(2) = 8$ and
$L^{\min}_{\mathrm{MM+}}(2) = 7$; the MM value agrees with Lemma~\ref{lem:snapshot} ($2n + k_{\mathrm{RS}} = 4 + 4$). For $n = 3$
an independent iterative-deepening search over the same exact states, with an admissible bound (cells off goal, missing
goal values, and $\lceil\log_2\rceil$ of the largest class of inputs not yet told apart), finds no MM+ kernel of length
$\le 14$ ($1.2 \times 10^9$ nodes, 35 minutes); so $L^{\min}_{\mathrm{MM+}}(3) = 15$. This is a search result without a
certificate; the same search finds the known 15-instruction kernels, a check that its pruning keeps the optimum. A shortest MM+ sort3 even exists in $I_0$.

\paragraph{Shorter is not faster here.} AlphaDev uses length as the cost because length and latency correlate for its
kernels. With memory operands this link breaks: in our measurements on the AMD EPYC 9654 the MM+ kernels have a longer
dependency-chain latency than AlphaDev's (median cycles per call, harness baseline subtracted: $n = 3$: 13.5 vs.\ 9.7;
$n = 4$: 15.4 vs.\ 10.4; $n = 5$, 40-instruction kernel: 24.2 vs.\ 16.7) and no better throughput, because each memory operand puts a load on the
critical path. (These timings come from a shared, loaded machine, so the absolute values are approximate.) So MM+ is a statement about AlphaDev's \emph{cost model}
(instruction count), not a faster sort.

\paragraph{Consequence for the main result.} AlphaDev's lengths 17, 28 and 42 are optimal in AlphaDev's own action space
(MM) for $n = 3, 4$ (and not beaten for $n = 5$), but not in the full x86 instruction set restricted to \texttt{mov},
\texttt{cmp} and \texttt{cmov}: the rule MM-1 is necessary for the claim. We state the main theorem for MM and report MM+ separately.


\section{Threats to validity: what is outside the model}\label{sec:threats}
\begin{itemize}
\item \textbf{Memory operands (MM+).} x86 allows \texttt{cmp r,m}, \texttt{cmp m,r}, \texttt{cmovcc r,m}. With them, sort4 and sort5
fit in 26 and 40 instructions (Section~\ref{sec:beyond}). So AlphaDev's kernels are optimal only under MM-1. Our snapshot
lemma does not hold for MM+ (Lemma L8: only $L \ge k_{\mathrm{RS}}$); an MM+ optimum needs its own encoding.
\item \textbf{Other instructions}: SIMD min/max, branches, arithmetic and masks, \texttt{setcc}, \texttt{xchg}, \texttt{lea},
\texttt{cmove/cmovne}, unsigned conditions (they read flags that are not a function of the signed order).
\item \textbf{Other costs}: latency, throughput, code size (AlphaDev reports that length and latency correlate for fixed sorts).
\item \textbf{Data}: signed int32. Unsigned or floating-point kernels use other opcodes.
\item \textbf{Registers}: our lower bounds allow unbounded registers (stronger); real x86 has at most 14 free GPRs.
\item \textbf{Trust base}: the CNF encoders (two independent ones, differential tests, Lean route \pending), \texttt{cake\_lpr},
the Lean kernel, the x86 semantics in Lean versus real hardware (differential tests).
\item \textbf{Relaxation}: sound, but it may admit programs that read garbage registers usefully; a found relaxed program is
re-checked in the strict model.
\end{itemize}

% spec, 2026-10-24. Section "Related work". Source notes: agents/spec/prior_work.md (primary sources in agents/spec/ext/).
\section{Related work}\label{sec:related}

\paragraph{AlphaDev.} See Section~\ref{sec:background}. AlphaDev's only optimality statement is the brute-force claim for sort3
(about $10^{32}$ programs, with pruning, no certificate, search space not stated)~\cite{alphadev}. The kernels went into
LLVM libc++ (review D118029, landed 2022-04-08 as commit \texttt{194d1965d2c8})~\cite{d118029}, where they were accepted on
benchmarks; the patch makes no optimality claim.

\paragraph{Ullrich and Hack.} The closest work is ``Synthesis of Sorting Kernels''~\cite{ullrich} (CGO 2025). They search
sorting kernels with A$^\ast$/Dijkstra over abstract states, with optimality-preserving heuristics and optional
non-optimality-preserving cuts. They enumerate all optimal kernels for $n = 3$ (length 11 without memory instructions,
5602 programs after symmetry reduction) and report that no kernel of length 19 exists for $n = 4$ (exhaustive, two weeks), a
``new tight lower bound''; for $n = 5$ they find kernels of length about 33 with a cut, and give no lower bound.
Table~\ref{tab:uh} compares their model with ours, point by point. Each point where their model is \emph{more restrictive}
than AlphaDev's (fewer programs) means that their lower bound excludes fewer programs and so does not imply a lower bound
for AlphaDev's setting.

\begin{table}[h]\centering\small
\begin{tabular}{p{2.6cm}p{5.2cm}p{5.2cm}}\toprule
 & Ullrich and Hack~\cite{ullrich} & this work (MM via RS) \\\midrule
instructions & \texttt{mov}, \texttt{cmp}, \texttt{cmovl}, \texttt{cmovg} ($I_0$) & $I_1$ (adds \texttt{cmovle}, \texttt{cmovge}), as AlphaDev; $I_0$ as a side result \\
memory, counting & none; registers only (add $2n$ to compare) & loads and stores counted; exact map $L = 2n + k$ (Lemma~\ref{lem:snapshot}) \\
outputs & fixed in $r_1..r_n$ at the end (RF) & any register at any step (RS) = AlphaDev's stores \\
scratch registers & a fixed number $m$; the paper does not state $m$; the artifact code uses one (\texttt{SWAPS = 1}) & unbounded ($R = n + k$ suffices, Lemma~L3) \\
inputs tested & the $n!$ permutations (no ties) & all weak orders (ties); permutation subsets only inside sound relaxations \\
initial state & registers hold the inputs; reads of uninitialised scratch registers are forbidden & strict: garbage registers and flags (reads of uninitialised registers allowed); lower bound via a sound relaxation \\
pruning rules & argued in prose & each rule with a paper proof (and Lean where marked) \\
evidence & exhaustive-search runs, no proof log & DRAT/LRAT certificates checked by \texttt{cake\_lpr}; verified checkers in Lean and Rocq; two independent search tools \\
results & $n=3$: 11; $n=4$: 20; $n=5$: none & $n=3$: 11 (certified); $n=4$: 20 (exhaustive search, certificate checked to $k \le 17$); $n=5$: 32 best known, $\le 31$ open \\\bottomrule
\end{tabular}
\caption{Ullrich and Hack's model versus ours. In their model a lower bound of 11 / 20 does not imply one for AlphaDev's
setting: $I_0 \subsetneq I_1$, RF is stricter than RS, one scratch register is stricter than unbounded registers, and
forbidding reads of uninitialised registers excludes programs that are correct on x86 (e.g.\ a 0-\texttt{mov} sort3 core with one
such read, found by our search).}
\label{tab:uh}
\end{table}

\noindent Their $n = 3$ and $n = 4$ values agree with ours. Their lower bound transfers to tie-correct kernels in $I_0$ with one
scratch register, fixed outputs and no reads of uninitialised registers (distinct-input correctness is weaker than tie correctness), but not to $I_1$, to more
scratch registers, or to snapshot outputs. They also report that SMT, CEGIS, SyGuS (cvc5), ILP, PDDL planners and STOKE do not
solve $n = 4$, and that STOKE fails already for $n = 3$.

\paragraph{Shorter kernels outside the model.} Neri~\cite{neri} gives a 15-instruction branchless sort3 and a 14-instruction
sort3 with a loop; both use \texttt{sbb}/\texttt{adc}, table lookups and indexed stores, so they are outside the
\texttt{mov}/\texttt{cmp}/\texttt{cmov} model. Tunney~\cite{tunney} explains the AlphaDev saving as the deletion of one
\texttt{mov} from three compare-and-swap blocks (her 17-instruction x86 kernel passes our judge; it needs \texttt{cmovge}) and
shows a 14-instruction ARM64 kernel with the three-operand \texttt{csel}, also outside the model. Mimicry~\cite{mimicry} computes
a shuffle index instead of swapping: a 9-instruction SIMD sort3, and scalar kernels with indexed stores (sort3 17, sort4 27
for signed values); these use instructions outside our model and make speed claims only. None of them claims or refutes
optimality in AlphaDev's model.

\paragraph{Superoptimization.} Massalin's superoptimizer~\cite{massalin}, STOKE~\cite{stoke}, Lens~\cite{lens} and
Souper~\cite{souper} find short or fast instruction sequences by enumeration, stochastic search or solvers; none of them reports
an optimal sorting kernel, and STOKE fails on sort3 in the experiments of~\cite{ullrich}.

\paragraph{Optimal sorting networks.} The comparator model is solved for small $n$ with SAT and computer-checked certificates:
9 inputs need 25 comparators~\cite{codish}, with a Coq-verified checker of the certificate~\cite{cruzfilipe}, and 11 and 12
inputs need 35 and 39 comparators, with an Isabelle-verified checker~\cite{harder}. These results count comparators, not
instructions, and the 0-1 principle holds there. For cmov code a comparator is not atomic: one \texttt{cmp} can feed
\texttt{cmov}s that do not form a full swap (the swap move), the instruction count is not $4 \times$ comparators, and the
0-1 principle fails (SI App.~F of~\cite{alphadev}). Our method follows the same pattern (SAT encoding, symmetry breaking,
verified certificate checking) for the instruction model.


\appendix
% Appendix: soundness proofs (theory, 2026-10-24). Source of truth: agents/theory/lemmas.md.
% Needs: amsmath, amsthm, amssymb. Macros below.
\newcommand{\lt}{\mathsf{lt}}\newcommand{\eqf}{\mathsf{eq}}\newcommand{\gt}{\mathsf{gt}}
\newcommand{\Wn}{\mathcal{W}_n}\newcommand{\srt}{\mathrm{sort}}
\newcommand{\RS}{\ensuremath{\mathrm{RS}}}\newcommand{\MM}{\ensuremath{\mathrm{MM}}}

\section{Soundness proofs}\label{app:proofs}

\subsection{Machine, semantics, output conventions}
Registers $\mathrm{Reg}=\{P_0,\dots,P_{n-1}\}\cup\{S_0,S_1,\dots\}$, a flag $F\in\{\lt,\eqf,\gt\}$.
Instructions (destination first): $\texttt{mov}\,D,S$; $\texttt{cmp}\,A,B$ sets $F$ to the outcome of comparing $A$ with $B$;
$\texttt{cmov}c\,D,S$ for $c\in\{l,g,le,ge\}$ sets $D:=S$ iff $F\in C(c)$ with
$C(l)=\{\lt\}$, $C(g)=\{\gt\}$, $C(le)=\{\lt,\eqf\}$, $C(ge)=\{\gt,\eqf\}$. $I_0$ uses $l,g$; $I_1$ uses all four.
Reads and writes: $W(\texttt{mov}\,D,S)=\{D\}$, $R(\texttt{mov}\,D,S)=\{S\}$; $W(\texttt{cmp}\,A,B)=\{F\}$, $R=\{A,B\}$;
$W(\texttt{cmov}c\,D,S)=\{D\}$, $R=\{D,S,F\}$. Only \texttt{mov} is a \emph{full definition} of its destination.

A run of $p=i_1\cdots i_k$ from a state $\sigma_0$ gives states $\sigma_0,\dots,\sigma_k$ (position $t$ = after $i_t$).
\emph{Strict} correctness quantifies over all inputs $x\in\mathbb{Z}_{32}^n$, all scratch values and all $F$.
The real flags $(ZF, SF\oplus OF)$ have four states; the four conditions read them as
$\lt$ iff $SF\neq OF$, $\eqf$ iff $SF=OF\wedge ZF$, $\gt$ iff $SF=OF\wedge\neg ZF$, so the three-valued $F$ is exact.
\emph{T-relaxed} correctness: inputs with $x_i<\mathrm{INT\_MAX}$, all $S_j=T:=\mathrm{INT\_MAX}$, $F=\eqf$.

An \RS{} program is a pair $(p,o)$ with $o(j)=(r_j,t_j)$, $0\le t_j\le k$; it is correct on $x$ if
$\sigma_{t_j}(r_j)=\srt_j(x)$ for all $j$. (RF: $o(j)=(P_j,k)$. RE: $t_j=k$.)

\subsection{Order types (L1)}
Let $\Wn$ be the weak orders on $\{0..n-1\}$ ($|\Wn|=1,3,13,75,541$ for $n=1..5$).
\emph{Abstract run} on $w\in\Wn$: registers hold labels in $\{0..n-1\}\cup\{T\}$, compared by $w$, with $T$ above every
index and $\mathrm{cmp}(T,T)=\eqf$.
\begin{lemma}[L1]
Let $\tau(x)=w$ and $x_i<\mathrm{INT\_MAX}$. In the T-relaxed run on $x$, at every position every register holds
$v_x(\ell)$, where $\ell$ is its label in the abstract run on $w$, $v_x(i)=x_i$, $v_x(T)=\mathrm{INT\_MAX}$; and the
flags agree. Hence $(p,o)$ is relaxed-correct on $x$ iff it is abstractly correct on $w$.
\end{lemma}
\begin{proof}
Induction on positions; the only nontrivial case is \texttt{cmp}: $v_x(\ell)<v_x(\ell')$ iff $\ell<_w\ell'$ because
$\tau(x)=w$ and every $x_i<\mathrm{INT\_MAX}$. For outputs: $v_x(\ell)=\srt_j(x)$ iff $\ell\neq T$ and $x_\ell=\srt_j(x)$,
a property of $w$ alone.
\end{proof}
\begin{corollary}[L1b]
For $W\subseteq\Wn$: if no length-$k$ program is abstractly correct on $W$, none is relaxed-correct.
\end{corollary}

\subsection{Relaxation (L2) and registers (L3)}
\begin{lemma}[L2] Strict correctness implies T-relaxed correctness (all output conventions). \end{lemma}
\begin{proof} The relaxed run is the strict run with scratch $=\mathrm{INT\_MAX}$ and the real flag state $ZF=1$, $SF=OF$. \end{proof}
\begin{lemma}[L3] If a relaxed-correct \RS{} program of length $k$ exists ($n\ge2$), one exists over $P_0..P_{n-1},S_0..S_{k-1}$. \end{lemma}
\begin{proof}
With no \texttt{cmp}, $F=\eqf$ throughout; $l,g$-cmovs are no-ops and $le,ge$-cmovs are \texttt{mov}s, so every label is
independent of $w$, and output $0$ fails on some $w$. Hence $\le k-1$ instructions write registers. Never-written scratch
registers hold $T$ forever: merge them into one. Rename the at most $k$ scratch names bijectively; all scratch registers
start equal, so the run is unchanged.
\end{proof}

\subsection{Memory model and the counting map (L5)}
An \MM{} program uses cells $m_0..m_{n-1}$ (holding $x$ at the start, required to hold $\srt(x)$ at the end) and
registers. Instructions: register instructions as above, \emph{loads} $\texttt{mov}\,r,[m_i]$, \emph{stores}
$\texttt{mov}\,[m_i],r$. Assumption (MM-1): no other instruction touches a cell.
\begin{lemma}[L5]
If an \MM{} program $q$ of length $L$ is correct (strict, or T-relaxed), then an \RS{} program of length $\le L-2n$ is
correct in the same semantics. Conversely an \RS{} program of length $k$ yields an \MM{} program of length $k+2n$.
\end{lemma}
\begin{proof}
\emph{(a) Each cell is loaded before it is first stored to.} Otherwise $x_i$ never reaches a register (MM-1), the final
memory does not depend on $x_i$, but $\srt(x)$ does.
\emph{(b) Each cell is stored to.} Otherwise $m_j$ ends as $x_j\neq\srt_j(x)$ for some $x$.

\emph{Segments.} The stores to $m_i$ split time into segments $s=0,1,\dots,e_i$; segment $0$ is before the first store and
has value $x_i$; segment $s\ge1$ starts at the $s$-th store. By (a) segment $0$ contains a load; by (b) $e_i\ge1$.
Let $G_{i,0}:=P_i$ and let $G_{i,s}$ ($s\ge1$) be fresh registers.

\emph{Live ranges.} A load or \texttt{mov} into $r$ at time $t$ starts a range of $r$ that ends just before the next
load or \texttt{mov} into $r$ (or at the end). A \texttt{cmov} into $r$ stays in the current range.

\emph{Translation $q\mapsto(p,o)$}, instruction by instruction, with a name map $\nu$ (initially identity on registers):
\begin{itemize}
\item register instruction: emit it with operands renamed by $\nu$;
\item load $\texttt{mov}\,r,[m_i]$ in segment $s$: if it is not the last load of the segment, emit $\texttt{mov}\,\nu'(r),G_{i,s}$
      with $\nu'(r)$ a fresh name, set $\nu(r):=\nu'(r)$; if it is the last load, emit nothing and set $\nu(r):=G_{i,s}$;
\item store $\texttt{mov}\,[m_i],r$ that opens segment $s$: if the segment has a load, emit $\texttt{mov}\,G_{i,s},\nu(r)$;
      otherwise emit nothing. If $s=e_i$, set $o(i):=(\nu(r),\text{current position})$.
\end{itemize}
\emph{Invariant.} After translating a prefix of $q$ (same semantics, corresponding initial states; fresh names start
as garbage, resp.\ $T$): (I1) for every register $r$, the value of $r$ in $q$ equals the value of $\nu(r)$ in $p$
whenever $r$'s current range is read later; (I2) for every cell $i$ in segment $s$ whose remaining loads are not
yet all done, $G_{i,s}$ in $p$ holds the segment value; (I3) the flags agree.
Each step keeps the invariant: register instructions act on corresponding values (I1, I3); a non-last load copies the segment
value by (I2); at the last load, $G_{i,s}$ holds the segment value and is not read again except through $\nu(r)$ (no later
load in this segment; later segments use other names), so (I1) holds for the new range; a store with a later load writes the
segment value to $G_{i,s}$ (I2). Two ranges never share a name: fresh names are fresh, and $G_{i,s}$ is renamed into at most
one range (the last load), after which nothing else reads or writes $G_{i,s}$. At the final store of $m_i$, (I1) gives
$\sigma(\nu(r))$ = the value stored, which is the final content of $m_i$, i.e.\ $\srt_i(x)$. So $(p,o)$ is correct.

\emph{Count.} Per cell: segment 0 loses its last load ($-1$); segment $e_i\ge1$ loses its store (no load) or its last
load ($-1$). Every other emitted instruction replaces one instruction of $q$. Hence $|p|\le L-2n$.

\emph{Converse.} Emit $\texttt{mov}\,P_i,[m_i]$ for all $i$, then $p$, inserting $\texttt{mov}\,[m_j],r_j$ at position $t_j$.
Cells are not read after the first $n$ instructions, so the stores are safe.
\end{proof}
\begin{remark} (i) The proof needs neither ``loads first / stores last'' nor AlphaDev's read-once/write-once rule.
(ii) Fixing both inputs and outputs (RF) is \emph{not} implied: \MM{} chooses load targets and store sources freely.
(iii) Without MM-1 (memory sources in \texttt{cmp}/\texttt{cmov}) a load can merge into a compare and (a) gives no
load per cell; the lemma then fails as stated.
\end{remark}

\subsection{Padding (L4)}
\begin{lemma}[L4] (a) Without symmetry rules, UNSAT at exact length $k$ implies UNSAT for all $k'\le k$
(append $\texttt{mov}\,P_0,P_0$). (b) With the rules of L6, use either the sweep ``UNSAT for every $k'\le K$'' or a
\texttt{nop} opcode allowed only as a suffix (rules apply to the non-\texttt{nop} prefix).
\end{lemma}
\begin{proof} (b) If a correct program of length $\le K$ exists, a shortest one has length $m^\ast\le K$ and by
Lemma~\ref{lem:master} a rule-satisfying one of length $m^\ast$ exists. \end{proof}

\subsection{Symmetry breaking and pruning (L6)}\label{lem:master}
Fix a total order $\prec$ on instructions (lexicographic in opcode, then operands, with a fixed register order) and
extend it lexicographically to programs of equal length (snapshot lists break ties).
Operations on relaxed-correct \RS{} programs of length $k$ (each preserves length and correctness):
$\rho_\pi$ permutes scratch names (all start as $T$);
$\iota_\pi$ permutes $P$ names, also in $o$ (the program then sorts $x\circ\pi$, and $\srt(x\circ\pi)=\srt(x)$; not valid in RF);
$\sigma_t$ swaps $i_t,i_{t+1}$ if they are \emph{independent}:
$W(i_t)\cap(R(i_{t+1})\cup W(i_{t+1}))=\emptyset=W(i_{t+1})\cap R(i_t)$, and moves a snapshot at position $t$ reading
$r$ to $t+1$ if $r\in W(i_t)$, to $t-1$ if $r\in W(i_{t+1})$;
$\phi_t$ replaces $\texttt{cmp}\,A,B$ at $t$ by $\texttt{cmp}\,B,A$ and swaps $l\leftrightarrow g$, $le\leftrightarrow ge$ in
the cmovs up to the next \texttt{cmp}.
Let $\mu$ replace every $le/ge$-cmov before the first \texttt{cmp} by \texttt{mov} (relaxed: $F=\eqf$ there).
\begin{lemma}[Master lemma]
Let $p$ be a shortest relaxed-correct program and $p^\ast=\min_\prec C(\mu(p))$, where $C$ is the closure under
$\rho,\iota,\sigma,\phi$. Then $p^\ast$ is shortest and correct, has no $le/ge$-cmov before the first \texttt{cmp}, and
$g(p^\ast)\not\prec p^\ast$ for every applicable operation $g$.
\end{lemma}
\begin{proof} $C(\mu(p))$ is finite and every member is correct of length $|p|$. No operation creates a pre-\texttt{cmp} cmov:
$\sigma$ never swaps a cmov with a \texttt{cmp} (both use $F$), $\rho,\iota$ keep opcodes, $\phi$ acts after a \texttt{cmp}. \end{proof}
Consequences (each is ``some $g$ would give a $\prec$-smaller program''):
R1 scratch names first occur in the order $S_0,S_1,\dots$ in the scan order of $\prec$;
R2 independent neighbours satisfy $i_t\prec i_{t+1}$ (identical instructions are never independent);
R5 every $\texttt{cmp}\,A,B$ has $A\prec B$;
R6 (RS/RE) $P$ names first occur in order;
R7 no $le/ge$-cmov before the first \texttt{cmp}.
Minimality rules (every shortest program satisfies them, since deleting the instruction keeps correctness):
S1 no $\texttt{mov}/\texttt{cmov}\,D,D$; S2 no dead write ($D$ not read --- as a source, as a cmov destination, by a
\texttt{cmp}, or by a snapshot --- before the next full definition of $D$ or the end); S3 every \texttt{cmp} is followed by a
cmov before the next \texttt{cmp}/end; S4 no $l/g$-cmov before the first \texttt{cmp}; S5 no $\texttt{cmp}\,A,A$.
Implied constraints (true for every correct program, e.g.\ ``a value needed by a later snapshot is held by some register'')
combine with all rules.
\begin{theorem}[Joint soundness]
If a relaxed-correct \RS{} program of length $\le K$ exists, then for $m^\ast$ = the least such length, a relaxed-correct
\RS{} program of length $m^\ast$ over $n+m^\ast$ registers satisfies R1, R2, R5, R6, R7, S1--S5 and every implied constraint.
\end{theorem}
\begin{proof} L3, then the master lemma applied to a shortest program over $n+m^\ast$ registers; $p^\ast$ is shortest, so S1--S5
hold. \end{proof}

\subsection{Comparison count (L7)}
\begin{lemma}[L7] A program correct on all $n!$ permutations (T-relaxed, any output convention, $I_0$ or $I_1$) has at least
$\lceil\log_2 n!\rceil$ \texttt{cmp} instructions. \end{lemma}
\begin{proof} Permutations with equal flag histories have equal labels in all registers (induction). So a \texttt{cmp} compares
one label pair on a whole class: equal labels or $T$ against anything give one outcome; distinct indices give $\lt$ or $\gt$.
Each \texttt{cmp} at most doubles the number of classes. At the end, snapshot $j$ holds label $\pi^{-1}(j)$, so the $n!$
permutations lie in distinct classes. \end{proof}

\subsection{Main theorem}
\begin{theorem}
Fix $n$, $I\in\{I_0,I_1\}$, MM-1, AlphaDev counting (every instruction, incl.\ loads and stores), unbounded registers.
If for every $k\le L-1-2n$ the encoding of ``a relaxed \RS{} program of length $k$ over $n+k$ registers, correct on
$W\subseteq\Wn$, satisfying a subset of the rules'' is unsatisfiable, then no \MM{} program of length $\le L-1$ sorts all
$x\in\mathbb{Z}_{32}^n$ for all garbage and flags.
\end{theorem}
\begin{proof} L2, L5, joint soundness, L1b. \end{proof}


\begin{thebibliography}{99}
% spec, 2026-10-24. \bibitem entries for background.tex, model_counting.tex, related.tex (keys alphadev, ullrich match theory's main.tex).
\bibitem{alphadev} D.~J. Mankowitz, A.~Michi, A.~Zhernov, M.~Gelmi, M.~Selvi, C.~Paduraru, E.~Leurent, et al.
Faster sorting algorithms discovered using deep reinforcement learning. \emph{Nature} 618(7964):257--263, 2023.
DOI 10.1038/s41586-023-06004-9. Supplementary Information: 41586\_2023\_6004\_MOESM1\_ESM.pdf.
\bibitem{alphadevrepo} Google DeepMind. AlphaDev repository, \texttt{sort\_functions\_test.cc}.
\url{https://github.com/google-deepmind/alphadev}, accessed 2026-10-24.
\bibitem{d118029} M.~Gelmi. Introduce branchless sorting functions for sort3, sort4 and sort5. LLVM review D118029,
\url{https://reviews.llvm.org/D118029}, landed 2022-04-08 (commit 194d1965d2c8).
\bibitem{ullrich} M.~Ullrich, S.~Hack. Synthesis of sorting kernels. In \emph{Proc.\ CGO 2025}, pp.~1--14.
DOI 10.1145/3696443.3708954. Artifact: Zenodo 10.5281/zenodo.14092980.
\bibitem{neri} C.~Neri. Shorter and faster than Sort3AlphaDev. arXiv:2307.14503, 2023.
\bibitem{tunney} J.~Tunney. Understanding DeepMind's sorting algorithm. \url{https://justine.lol/sorting/}, 2023-06-12
(read via web.archive.org).
\bibitem{mimicry} Mimicry. Faster sorting beyond DeepMind's AlphaDev.
\url{https://www.mimicry.ai/faster-sorting-beyond-deepminds-alphadev}, 2023 (read via web.archive.org).
\bibitem{massalin} H.~Massalin. Superoptimizer: a look at the smallest program. In \emph{Proc.\ ASPLOS 1987}, pp.~122--126.
\bibitem{stoke} E.~Schkufza, R.~Sharma, A.~Aiken. Stochastic superoptimization. In \emph{Proc.\ ASPLOS 2013}, pp.~305--316.
\bibitem{lens} P.~M. Phothilimthana, A.~Thakur, R.~Bodik, D.~Dhurjati. Scaling up superoptimization. In
\emph{Proc.\ ASPLOS 2016}. DOI 10.1145/2872362.2872387.
\bibitem{souper} R.~Sasnauskas, Y.~Chen, P.~Collingbourne, J.~Ketema, G.~Lup, J.~Taneja, J.~Regehr. Souper: a synthesizing
superoptimizer. arXiv:1711.04422, 2017.
\bibitem{codish} M.~Codish, L.~Cruz-Filipe, M.~Frank, P.~Schneider-Kamp. Sorting nine inputs requires twenty-five
comparisons. \emph{J.\ Comput.\ Syst.\ Sci.} 82(3):551--563, 2016.
\bibitem{cruzfilipe} L.~Cruz-Filipe, K.~S. Larsen, P.~Schneider-Kamp. Formally proving size optimality of sorting networks.
\emph{J.\ Autom.\ Reasoning} 59(4):425--454, 2017.
\bibitem{harder} J.~Harder. An answer to the Bose-Nelson sorting problem for 11 and 12 channels. arXiv:2012.04400, 2020.

\end{thebibliography}
\end{document}
